\chapter{Testing and Early Validation}

\section{Testing Strategy}

We test at four levels and add manual end-to-end checks against the running system. The aim is that every claim in this report maps to a test, a measurement or a screenshot, and that a failure is reported instead of hidden.

\begin{enumerate}
\item Unit tests of backend services with the language model replaced by a fake, so classification, routing, ranking and citation rules are tested deterministically.
\item API tests that call the routes through FastAPI's test client, including the failure codes (404, 409, 422, 503).
\item Repository contract tests that run the same graph operations against the in-memory store and a real Neo4j service.
\item Component tests of the React pages with the API mocked, covering every response shape and every error state.
\item Smoke and end-to-end checks. A Docker Compose job in CI starts the whole stack and verifies health, the loaded corpus and the maintenance APIs; the acceptance tests T01--T16 are run by hand against a running system.
\end{enumerate}

For the frontend work added late in the iteration we also applied mutation checks: after writing a test, we deliberately broke the code it protects (for example, removed a size limit or reverted a label) and confirmed that the test failed, then restored the code.

\section{Test Levels and Techniques}

Table~\ref{tab:tests-by-area} summarises the automated tests by area at the evaluated commit.

\begin{table}[H]
\centering
\caption{Automated tests by area (commit \texttt{36a98a3})}
\label{tab:tests-by-area}
\small
\begin{tabularx}{\textwidth}{@{}>{\ttfamily\raggedright\arraybackslash}p{0.42\textwidth} c L@{}}
\toprule
\normalfont File & Tests & What it covers \\
\midrule
\multicolumn{3}{@{}l@{}}{\textbf{Backend (PyTest)} --- 194 tests, of which 145 need no database} \\
\midrule
test\_conflict\_engine.py & 32 & Candidate pairing, classification, confidence, routing \\
test\_diagnostic\_retrieval.py & 26 & Bounded retrieval, fleet scoring, ranking tie-break \\
test\_graph\_maintenance.py & 49 & Schema, ingest links, Neo4j/memory contract \\
test\_api\_maintenance.py & 27 & Assets, history, records, ingest, review APIs \\
test\_maintenance\_agent.py & 10 & CSV parsing, skip categories, idempotent ingest \\
test\_conflict\_cascade.py & 11 & Deterministic screen, voting, cascade budget \\
test\_citation\_validation.py & 6 & Identifier shape, graph lookup, unverified marking \\
test\_api\_copilot.py & 10 & Multi-turn state, deltas, 503 empty explanation \\
test\_backfill\_conflict\_scores.py & 6 & Resume, budget stop, bounded retries \\
test\_api\_diagnostics.py & 4 & Diagnose contract and LLM gate \\
test\_ops.py & 3 & Version and ops log endpoints \\
test\_conflict\_analysis\_migration.py & 3 & Legacy-placeholder migration \\
test\_regression\_software.py & 7 & Frozen software-domain adapter isolation \\
\midrule
\multicolumn{3}{@{}l@{}}{\textbf{Frontend (Vitest)} --- 171 tests in 15 files} \\
\midrule
Diagnostics.test.jsx & 29 & Diagnose states, citations, evidence links \\
ServiceRecords.test.jsx & 22 & Fleet list, filters, pagination, deep links \\
ReviewQueue.test.jsx & 18 & Accept/reject confirmation, 409, required note \\
Graph.test.jsx & 18 & Graph selection, counts, edge styling \\
MaintenanceIngest.test.jsx & 17 & Upload confirmation, counts, error states \\
AssetHistory.test.jsx & 12 & Chronology, corrected-by banner, deep links \\
graphScope.test.js & 10 & Scope filtering helpers \\
diagnosticsUtils.test.js & 24 & Ranking/citation display helpers \\
DiagnosticsCopilot.test.jsx & 4 & Multi-turn UI state \\
Settings / Sidebar / DomainContext / dummy & 13 & Navigation and system status \\
Assets.test.jsx & 3 & Asset list and search \\
ingestClient.test.js & 1 & Upload client helper \\
\bottomrule
\end{tabularx}
\end{table}

\section{Test Cases}

\subsection{End-to-end acceptance tests}

The team defined T01--T16 before the iteration (owner in brackets). Table~\ref{tab:acceptance} gives each result and its evidence.

\begin{table}[H]
\centering
\caption{Acceptance tests and results}
\label{tab:acceptance}
\small
\begin{tabularx}{\textwidth}{@{}c >{\raggedright\arraybackslash}p{0.26\textwidth} >{\raggedright\arraybackslash}p{0.10\textwidth} L@{}}
\toprule
ID & Scenario & Result & Evidence and caveat \\
\midrule
T01 & Application opens the maintenance workspace & Pass & Route smoke log; Sidebar tests. Browser console was not audited systematically \\
T02 & API lists loaded assets & Pass & \texttt{test\_api\_maintenance}; staging lists 118 aircraft \\
T03 & Selected asset history is displayed & Pass & AssetHistory tests; Figure~6.7 \\
T04 & Diagnosis returns graph-backed evidence and ranking & Partly & Full result with a live model locally (Figure~6.2); on staging fresh diagnoses returned no explanation text on 16~Oct \\
T05 & Explanation cites a real record that opens & Pass & \texttt{test\_citation\_validation}; Figure~6.3 \\
T06 & Unknown asset request & Pass & Controlled 404, ``Aircraft not found'', no invented history \\
T07 & Model unavailable & Pass & HTTP 503 with a reason; Figure~6.5 \\
T08 & Empty or ambiguous evidence & Pass & Explicit ``no relevant history''; ranking withheld when there is no evidence \\
T09 & Review queue shows a known case & Pass & Figure~6.8 \\
T10 & Review action on a test case & Pass & ReviewQueue and API tests; accept/reject not exercised live on shared data \\
T11 & Frontend tests and build & Pass & 171 tests pass; production build passes \\
T12 & Backend tests & Pass & 194 tests pass in CI \\
T13 & Pull-request CI runs and passes & Pass & Run \texttt{39068741787} on \texttt{36a98a3} succeeded \\
T14 & Real graph health & Pass & Staging reports Neo4j, live model, 2{,}079 records (18~Oct); Compose smoke \\
T15 & Setup on a clean machine & Partly & CI builds from a clean checkout; nobody outside the team has followed the guide \\
T16 & Legacy screen isolation & Pass & Sidebar and DomainContext tests; legacy URLs redirect \\
\bottomrule
\end{tabularx}
\end{table}


Appendix~C lists every automated test by name.

\section{Test Results}

At commit \texttt{36a98a3} the CI run (id \texttt{39068741787}) reported: backend 194 passed; frontend 171 passed in 15 files; the Docker Compose real-graph smoke test passed; the staging deployment that followed succeeded. Code coverage has not been measured (NFR-04), so no coverage figure is claimed.

\section{Acceptance Testing / Validation}

\subsection{One-aircraft analysis run}

To show scores, corrections and a review queue on real records, and to test the conflict engine on a bounded input, we analysed one aircraft, 800NN (13 records, three FAA-recorded correction pairs), on a throwaway Neo4j graph. Shared staging was not touched. The records were loaded through the new upload route and analysed in the original (legacy) mode with \texttt{openai/gpt-oss-120b} and the 0.50 threshold, with an empty classification cache so every result came from a real model call.

\begin{table}[!htbp]
\centering
\caption{Single-aircraft analysis results (800NN)}
\label{tab:single-aircraft}
\small
\begin{tabularx}{\textwidth}{@{}L r@{}}
\toprule
Measure & Result \\
\midrule
Records processed & 13 \\
Model classifications / screened without a model & 11 / 2 \\
Successful model requests & 55 \\
Tokens used & 29{,}386 \\
Rate-limit (HTTP 429) responses, all retried successfully & 27 \\
SUPERSEDES edges created & 4 \\
Status: no\_action / auto\_accepted / needs\_review & 7 / 4 / 2 \\
FAA-recorded pairs found & 3 of 3 \\
Edges not in the FAA list & 1 \\
\bottomrule
\end{tabularx}
\end{table}

The FAA ground-truth file lists only corrections the FAA reporters themselves recorded, so the one extra edge is unverified, not shown to be wrong. We do not compute a precision from one aircraft. Two records were sent to human review. These figures describe this aircraft only: about 2{,}260 tokens and 4.2 requests per record must not be extrapolated to the full corpus. The method and results are in \path{docs/mid1/23_single_aircraft_analysis_evidence.md}.

\subsection{Earlier offline evaluation of conflict detection}

An earlier evaluation measured the conflict engine on 350 pairs (141 FAA-recorded positives and 209 sampled negatives, one evaluation per superseding record), with \texttt{openai/gpt-oss-120b} at threshold 0.50:

\begin{table}[!htbp]
\centering
\caption{Offline conflict-detection evaluation}
\label{tab:offline-eval}
\small
\begin{tabularx}{\textwidth}{@{}L c c c c@{}}
\toprule
System & Precision & Recall & F1 & False-conflict rate \\
\midrule
Baseline (most recent prior) & 0.380 & 0.908 & 0.536 & 1.000 \\
Engine (threshold 0.50) & 0.903 & 0.660 & 0.762 & 0.048 \\
\bottomrule
\end{tabularx}
\end{table}

The baseline reaches its recall by asserting a supersession on every pair it sees, which is why its false-conflict rate is 1.000. Measured at pair level over the 141 FAA-recorded pairs, precision is 0.788 and recall 0.660; the two precision figures answer different questions and must not be mixed.

\textbf{How far this evidence goes.} It is historical: it comes from saved result files of the older run, with a model that is not the one now configured on staging. It measures conflict detection and is not a measure of diagnostic accuracy. The full-corpus counts published with it (121 detected, 117 auto-resolved, 4 for review, 1{,}617 without candidates) cannot be reproduced from the current code and are not used here.

\section{Edge and Failure Cases}

Table~\ref{tab:edge-cases} records the edge and failure cases exercised in this iteration and the behaviour under test.

\begin{table}[H]
\centering
\caption{Edge and failure cases, and the behaviour tested}
\label{tab:edge-cases}
\small
\begin{tabularx}{\textwidth}{@{}>{\raggedright\arraybackslash}p{0.30\textwidth} L@{}}
\toprule
Case & Behaviour \\
\midrule
Unknown aircraft & 404 and ``Aircraft not found'', never an empty history \\
Aircraft looked up by tail number & History route takes the UUID; a tail number returns 404 \\
No live model & HTTP 503 with a reason; Copilot returns an empty reply but keeps ranking and evidence \\
Empty model output & Explicit 503, never a blank success \\
Invented record identifier & Marked \texttt{[UNVERIFIED]} and left out of the valid list \\
Rate limit & Bounded retries; a daily budget stops runs cleanly \\
Superseded record in the Copilot & Visible as evidence, excluded from recommendations \\
Legacy confidence 1.0 & Reported as unscored with null confidence \\
Review action on a non-reviewable record & HTTP 409; the queue is unchanged \\
Reject without a note & Rejected by validation (HTTP 422) \\
Oversize or non-CSV upload & Rejected in the browser before anything is sent \\
Re-upload of the same CSV & Every record counted as already stored; nothing changes \\
Unreadable CSV & Backend may raise HTTP 500; interface wording discloses this (open issue) \\
Browser storage denied or corrupt & Copilot remains usable; a warning is shown \\
\bottomrule
\end{tabularx}
\end{table}


\section{Defects and Corrective Actions}

Table~\ref{tab:defects} lists defects found during the iteration and their resolution state.

\begin{table}[H]
\centering
\caption{Defects found and their resolution}
\label{tab:defects}
\small
\begin{tabularx}{\textwidth}{@{}c L L >{\raggedright\arraybackslash}p{0.10\textwidth}@{}}
\toprule
\# & Defect & Corrective action & State \\
\midrule
1 & Frontend CI failed: review-queue test mock lacked \texttt{getRecord} & Mock fixed; storage shim added & Fixed \\
2 & Review-queue history links used tail numbers and led to 404 & Resolve aircraft UUID; hide link if unresolved & Fixed \\
3 & Deep link did not scroll while loading skeleton was shown & Scroll effect depends on loading state & Fixed \\
4 & Service Records issued 1+118 requests on load & Fleet-wide paginated endpoint; two page-initiated calls & Fixed \\
5 & Assets list was slow (one history query per aircraft) & Single aggregation query; $\approx$1.0\,s on staging & Fixed \\
6 & Placeholder confidence 1.0 shown as a real result & Honest analysis state; placeholders reported as unscored & Fixed \\
7 & Graph requested 200 nodes yet claimed ``Capped at 2000'' & Request 3{,}000; truthful counts; warning when cap is hit & Fixed \\
8 & CSV upload client posted to a path-only route & Posts to the upload route; new upload panel & Fixed \\
9 & Small graphs hid edges behind large nodes & Spacing and node size changed & Fixed \\
10 & Staging fresh diagnoses returned no explanation text & Diagnose-prompt bug; API reports empty output as 503. Fixed after \texttt{36a98a3} (Appendix~D) & Closed after eval \\
11 & Unreadable CSV returns unhandled HTTP 500 & Interface wording only; backend change still open & Open \\
12 & One staging record has no ABOUT link & Re-running ingest repairs it; not applied without approval & Open \\
\bottomrule
\end{tabularx}
\end{table}


\section{Current Validation Status}

\textbf{Validated:} the code paths covered by 365 automated tests; the end-to-end slice on real records; the failure behaviour in Table~\ref{tab:edge-cases}; the engine on one aircraft against the FAA's own labels (3 of 3).

\textbf{Not validated:} diagnostic accuracy of the ranking; retrieval quality against baselines; the engine on the full corpus or with the staging model; usability with real technicians; load performance; code coverage; setup by an outside person. These are planned in Chapter~9.
